\begin{frame}{Introducción}
  \begin{itemize}
    \item El STM no es un microscopio óptico convencional; no usa lentes ni
      luz.
    \item Se basa en el \textbf{efecto túnel cuántico}.
    \item Su invención en 1981 marcó el inicio de la era de la nanotecnología.
    \item Desarrollado por \textbf{Gerd Binnig} y \textbf{Heinrich Rohrer} en
      IBM Zúrich.
  \end{itemize}
\end{frame}

\begin{frame}{1927: Friedrich Hund y el Origen de la Idea}
  \begin{block}{La Paradoja de los Isómeros}
    Friedrich Hund, mientras estudiaba la estructura de moléculas como el
    amoníaco ($\ce{NH3}$), observó algo inexplicable para la física clásica.
  \end{block}
  \begin{itemize}
    \item \textbf{Barreras de Potencial:} El átomo de nitrógeno oscilaba entre
      dos posiciones de equilibrio a pesar de no tener energía suficiente para
      saltar la barrera de repulsión de los hidrógenos.
    \item \textbf{Primer concepto de Túnel:} Hund propuso que, debido a la
      naturaleza ondulatoria de la materia, existe una probabilidad no nula de
      que una partícula atraviese una barrera energética.
    \item \textbf{Importancia:} Fue la primera vez que se aplicó el concepto de
      ``túnel'' para explicar la estabilidad y simetría molecular.
  \end{itemize}
\end{frame}

\begin{frame}{1928: La Formalización del Efecto Túnel}
  El año 1928 fue crucial para que el túnel pasara de ser una conjetura a una
  ley física aceptada:
  \begin{itemize}
    \item \textbf{George Gamow:} Explicó la desintegración alfa de los núcleos
      atómicos. Sin el efecto túnel, el núcleo sería una ``prisión'' eterna para
      las partículas alfa.
    \item \textbf{Oppenheimer y Fowler:} Aplicaron de forma independiente la
      mecánica cuántica a la emisión de electrones desde metales bajo campos
      eléctricos intensos (Emisión de Campo).
    \item \textbf{Conclusión Teórica:} Se entendió que la corriente de túnel es
      \textbf{exponencialmente dependiente} del ancho de la barrera. Esta
      sensibilidad es el corazón del futuro STM.
  \end{itemize}
\end{frame}

\begin{frame}{1930 - 1960: El Túnel en la Materia Condensada}
  Durante estas décadas, el efecto túnel dejó de ser un fenómeno nuclear para
  convertirse en una herramienta de estado sólido:
  \begin{itemize}
    \item \textbf{Aproximación JWKB:} Se perfeccionaron los métodos matemáticos
      para calcular probabilidades de túnel en barreras de formas arbitrarias.
    \item \textbf{Leo Esaki (1958):} Demostró el efecto túnel en semiconductores
      (Diodo Esaki), probando que se podía generar una corriente eléctrica
      puramente por tunelización a nivel industrial.
    \item \textbf{Ivar Giaever (1960):} Realizó experimentos de túnel en
      superconductores usando capas de óxido, demostrando que la corriente túnel
      permitía medir la ``densidad de estados'' (qué niveles de energía están
      ocupados por electrones).
  \end{itemize}
\end{frame}

\begin{frame}{1971: Russell Young y el Topografiner}
  El antecedente directo del STM fue el ``Topografiner'', desarrollado en el
  National Bureau of Standards (EE. UU.):
  \begin{itemize}
    \item \textbf{La Propuesta:} Russell Young intentó usar una punta metálica
      afilada para rastrear una superficie mediante la corriente de emisión de
      campo.
    \item \textbf{Aportes Clave:}
      \begin{itemize}
        \item Uso de actuadores piezoeléctricos para movimiento ultra-fino.
        \item Un sistema de control de retroalimentación (feedback loop).
      \end{itemize}
    \item \textbf{El Fracaso:} No lograba resolución atómica debido a que
      operaba en régimen de ``emisión de campo'' (distancia grande) y sufría de
      vibraciones mecánicas insuperables para la época.
  \end{itemize}
\end{frame}

\begin{frame}{1978 - 1981: El Desafío de IBM Zúrich}
  Gerd Binnig y Heinrich Rohrer se propusieron lo que muchos consideraban
  imposible: medir corrientes de túnel a través del vacío con una precisión de
  picómetros.
  \begin{itemize}
    \item \textbf{El Problema del Ruido:} Para ver átomos, el instrumento debía
      ser inmune a las vibraciones de un edificio o incluso a los pasos de
      alguien en el pasillo.
    \item \textbf{Ingeniería de Aislamiento:}
      \begin{itemize}
        \item Construyeron un sistema de levitación magnética sobre un cuenco de
          plomo superconductor para aislar el microscopio.
        \item Diseñaron el ``Piojo'' (louse), un motor piezoeléctrico capaz de
          dar ``pasos'' microscópicos para acercar la punta a la muestra sin
          chocar.
      \end{itemize}
  \end{itemize}
\end{frame}

\begin{frame}{1981: El Nacimiento del STM}
  \begin{itemize}
    \item \textbf{16 de Marzo de 1981:} Se registra por primera vez una
      corriente de túnel estable y controlada a través de un espacio de vacío
      real.
    \item \textbf{Primera Imagen:} Se analizó una superficie de oro ($\ce{Au}$),
      donde se observaron escalones atómicos.
    \item \textbf{Escepticismo:} La comunidad científica fue inicialmente
      incrédula; muchos pensaban que los resultados eran ruido electrónico o
      artefactos de la punta.
  \end{itemize}
\end{frame}

\begin{frame}{1982: La Reconstrucción del Silicio (7x7)}
  El STM se validó definitivamente al resolver el ``enigma del silicio'':
  \begin{itemize}
    \item La superficie del Silicio (111) presentaba una estructura extraña
      llamada $7 \times 7$ que nadie entendía desde 1959.
    \item \textbf{El Resultado:} Binnig y Rohrer publicaron una imagen donde se
      veían claramente los 12 ``adatoms'' de la celda unidad.
    \item \textbf{Impacto:} Fue la primera vez que el ojo humano ``vio'' la
      disposición real de los átomos en un cristal. El debate sobre si el STM
      funcionaba se terminó instantáneamente.
  \end{itemize}
\end{frame}

\begin{frame}{1986: El Premio Nobel y el Legado}
  \begin{itemize}
    \item \textbf{Reconocimiento:} Solo cinco años después de su invención,
      Binnig y Rohrer recibieron el Premio Nobel de Física.
    \item \textbf{Revolución Instrumental:} Del STM derivaron otros
      microscopios, como el de Fuerza Atómica (AFM), permitiendo estudiar
      materiales no conductores.
    \item \textbf{Manipulación Atómica (1990):} Don Eigler (IBM) usó un STM para
      mover 35 átomos de Xenón y escribir ``IBM'', demostrando que ahora podíamos
      no solo ver, sino construir a escala atómica.
  \end{itemize}
\end{frame}

